\chapter{系统需求分析}
\label{chap:requirements}

建筑幕墙韧性评估原本较多依赖人工整理和工程经验。将其转化为软件系统时，不能只从页面功能出发，还需要把现场工作中反复出现的信息拆分为可保存、可流转和可追溯的业务对象。本文将系统中的核心对象划分为项目、图像组、图像、检测主任务、检测子任务、图像总结、人工复核和项目报告。项目用于保存建筑名称、位置、建设年份、既有问题和评估目标等背景信息；图像组对应建筑立面、楼层或采集区域；图像记录原图、尺寸、文件大小和上传状态；检测主任务以单张图像为单位，串联分类路由、五类检测子任务和图像总结；人工复核记录工程人员对检测结果的补充判断；项目报告则面向后续运维决策形成最终输出。

从用户操作过程看，本系统也不是一组可以任意跳转的独立功能，而是围绕评估任务逐步推进的阶段化流程。用户通常需要先维护项目基础信息，再完成图像资产构建，随后启动 Agent 智能检测；检测结果经过人工复核确认后，系统才进入评估报告生成阶段。各阶段之间存在明确的前后依赖，例如没有图像资产就无法启动检测，检测结果未完成也不适合进入报告生成。基于上述对象模型和流程模型，本章进一步分析角色责任、典型用例、功能需求、非功能需求和相关约束条件。

\section{业务闭环中的角色与责任}
\label{sec:requirements-role-and-responsibility-in-business-closed-loop}

如\cref{fig:requirements-role-responsibility} 所示，系统中的主要角色可以分为使用者、业务服务、Agent 服务和检测模型服务。使用者关注的是项目创建、图像上传、结果查看、人工复核和报告确认；业务服务关注的是鉴权、数据持久化、状态推进和结果发布；Agent 服务关注的是语义路由、图像总结和项目报告生成；检测模型服务关注的是五类图像级原子检测能力的执行。角色划分的核心不是简单列出模块，而是明确“谁可以改变业务状态、谁只能提供判断结果”。这与\cref{sec:foundation-agent-engineering-controllability-problem}关于 Agent 工程化可控性问题的讨论直接对应。

从业务对象角度看，图像不是孤立输入，而是项目资产的一部分；检测任务也不是一次函数调用，而是围绕图像产生的业务记录。这样建模可以支持两个重要能力：第一，用户可以在任何时间回到项目查看历史检测结果；第二，系统可以在报告阶段追溯每个结论来自哪张图像、哪个检测任务和哪个区域。

\begin{figure}[!htb]
  \centering
  \begin{tikzpicture}[
    font=\footnotesize,
    >=stealth,
    role/.style={draw, rounded corners=2pt, align=center, minimum width=2.6cm, minimum height=0.95cm, line width=0.5pt, inner sep=2pt},
    service/.style={draw, rounded corners=2pt, align=center, minimum width=2.75cm, minimum height=0.95cm, line width=0.5pt, inner sep=2pt},
    wide/.style={draw, rounded corners=2pt, align=center, minimum width=8.1cm, minimum height=0.82cm, line width=0.5pt, inner sep=2pt},
    arrow/.style={->, line width=0.55pt},
    note/.style={align=center, font=\footnotesize},
    edge label/.style={font=\tiny, inner sep=1pt}
  ]
    \node[role, fill=blue!8] (user) at (0, 3) {使用者\\项目创建、图像上传\\结果查看与人工复核};
    \node[service, fill=cyan!8] (api) at (3.85, 3) {Core API\\HTTP 接入\\鉴权与 WebSocket};
    \node[service, fill=green!8, minimum width=3.1cm, minimum height=1.2cm] (biz) at (7.75, 3) {Core BIZ\\业务状态机\\任务创建与事务落库};

    \node[service, fill=orange!10] (classification) at (12.0, 5.25) {Classification\\分类 Agent\\输出检测任务代码};
    \node[service, fill=red!7] (reasoning) at (12.0, 3) {Reasoning\\检测模型服务\\执行原子检测};
    \node[service, fill=purple!8] (summary) at (12.0, 0.75) {Summary\\总结 Agent\\图像级总结与项目报告};

    \node[wide, fill=gray!10, minimum width=9.1cm] (infra) at (6.0, -0.65) {数据与事件底座：MySQL 保存结构化状态 \quad MinIO 保存图像产物 \quad RocketMQ 发布项目事件};
    \node[note] (principle) at (6.0, 6.35) {控制原则：Agent 与模型服务只回传结构化结果，Core BIZ 统一执行业务状态推进};

    \draw[arrow] ([yshift=0.16cm]user.east) -- node[edge label, above=2pt] {操作请求} ([yshift=0.16cm]api.west);
    \draw[arrow] ([yshift=-0.16cm]api.west) -- node[edge label, below=2pt] {状态展示} ([yshift=-0.16cm]user.east);
    \draw[arrow] ([yshift=0.16cm]api.east) -- node[edge label, above=2pt] {RPC 调用} ([yshift=0.16cm]biz.west);
    \draw[arrow] ([yshift=-0.16cm]biz.west) -- node[edge label, below=2pt] {查询结果} ([yshift=-0.16cm]api.east);

    \draw[arrow] ([xshift=-0.16cm]biz.north east) -- node[edge label, pos=0.55, above left=1pt] {图像路由任务} ([yshift=0.16cm]classification.south west);
    \draw[arrow] ([xshift=0.16cm]classification.south west) -- node[edge label, pos=0.55, below right=1pt] {任务代码} ([yshift=-0.16cm]biz.north east);

    \draw[arrow] ([yshift=0.16cm]biz.east) -- node[edge label, above=2pt] {检测子任务} ([yshift=0.16cm]reasoning.west);
    \draw[arrow] ([yshift=-0.16cm]reasoning.west) -- node[edge label, below=2pt] {指标与产物} ([yshift=-0.16cm]biz.east);

    \draw[arrow] ([xshift=0.16cm]summary.north west) -- node[edge label, pos=0.55, above right=1pt] {总结与报告} ([yshift=0.16cm]biz.south east);
    \draw[arrow] ([xshift=-0.16cm]biz.south east) -- node[edge label, pos=0.55, below left=1pt] {总结/报告任务} ([yshift=-0.16cm]summary.north west);

    \draw[arrow] ([xshift=-0.16cm]biz.south) -- node[edge label, left=2pt] {状态写入} ([xshift=-0.16cm]infra.north);
    \draw[arrow] ([xshift=0.16cm]infra.north) -- node[edge label, right=2pt] {文件与事件} ([xshift=0.16cm]biz.south);
  \end{tikzpicture}
  \caption{业务闭环中的角色与责任关系}\label{fig:requirements-role-responsibility}
\end{figure}

从用户流程角度看，五个阶段并不是并列入口，而是逐步补齐评估材料的过程。项目基础信息提供报告背景，图像资产提供检测对象；检测完成以后，人工复核才有具体依据，复核意见也会继续进入报告上下文。因此，系统需要设置项目进度控制，避免用户在数据尚未完整时提前进入后续阶段。与此同时，检测和报告生成都可能持续较长时间，若让用户一直等待同步请求返回，使用体验会明显下降，所以这类任务更适合以异步方式执行。

在角色边界上，各服务只负责自身范围内的判断或产物。分类 Agent 负责给出需要执行的检测任务建议，数据库中的检测子任务仍由业务服务创建；Reasoning 服务负责执行模型并生成检测产物，检测主任务是否完成则应由业务服务根据子任务状态统一判断；Summary 服务负责生成图像总结和项目报告，但项目阶段推进仍需要经过业务状态机；前端可以展示进度和结果，却不应根据局部页面状态自行推断数据库中的最终状态。

\begin{table}[!htb]
  \caption{业务角色与责任边界}\label{tab:role-boundary}
  \centering
  \begin{tabular}{@{}p{0.11\linewidth}p{0.4\linewidth}p{0.44\linewidth}@{}} \toprule
    \textbf{角色} & \textbf{负责内容} & \textbf{不应负责内容} \\ \midrule
    用户界面 & 用户交互、状态展示、结果查看、复核输入 & 直接修改任务状态或绕过后端判断项目是否完成 \\
    Core API & HTTP 接入、登陆鉴权、WebSocket 推送 & 复杂数据库事务和模型调度规则 \\
    Core BIZ & 业务状态机、任务创建、事务落库、事件发布 & 直接执行深度学习模型或自由生成报告文本 \\
    Activity 服务 & Agent 调用、模型推理、能力执行和回调 & 直接操作核心业务数据库 \\
    Agent 平台 & 语义路由、图像级总结、项目报告生成 & 无约束地改变业务流程或访问内部状态 \\ \bottomrule
  \end{tabular}
\end{table}

\cref{tab:role-boundary} 表明，本文系统的复杂性并不只是功能数量多，而是每个功能都需要明确归属。尤其在 Agent 引入后，必须避免让 Agent 获得过大的业务权限。本系统采用“Agent 输出建议，BIZ 执行业务”的方式，将核心状态推进集中在 Core BIZ 服务中，使智能能力和业务控制相互分离。这一设计虽然增加了 BIZ 的复杂度，但能够将 Agent 的语义能力和工程系统的确定性控制分离开来。

\section{典型用例分析}
\label{sec:requirements-typical-use-case-analysis}

为了进一步明确系统边界，可以将用户使用过程抽象为\cref{fig:requirements-core-use-case} 所示的系统核心 UML 用例图。土木工程评估用户是系统的主要参与者，直接触发登录鉴权、项目维护、图像资产管理、Agent 智能检测、结果复核和报告生成等用例。其中，启动 Agent 智能检测用例包含自动选择检测任务、执行原子检测能力和生成单图检测总结等内部用例；生成与查看评估报告用例包含整合项目报告上下文这一内部用例；失败重试与状态恢复作为扩展用例，用于处理检测或报告生成过程中的异常情况。

\begin{figure}[!htb]
  \centering
  \begin{tikzpicture}[
    font=\footnotesize,
    >=stealth,
    usecase/.style={draw, ellipse, fill=white, align=center, minimum width=2.45cm, minimum height=0.72cm, line width=0.5pt, inner sep=1pt},
    association/.style={-, line width=0.5pt},
    include/.style={->, dashed, line width=0.45pt},
    uml label/.style={font=\tiny, fill=white, inner sep=1pt}
  ]
    \draw[line width=0.5pt, rounded corners=2pt] (1.2, -4.4) rectangle (12.3, 4.05);
    \node[anchor=west] at (1.45, 3.75) {建筑幕墙韧性评估软件系统};

    \draw[line width=0.5pt] (0, 1.0) circle (0.16);
    \draw[line width=0.5pt] (0, 0.84) -- (0, 0.2);
    \draw[line width=0.5pt] (-0.34, 0.62) -- (0.34, 0.62);
    \draw[line width=0.5pt] (0, 0.2) -- (-0.3, -0.32);
    \draw[line width=0.5pt] (0, 0.2) -- (0.3, -0.32);
    \node[align=center] (user) at (0, -0.78) {土木工程\\评估用户};

    \node[usecase] (login) at (3.75, 2.85) {登录与\\身份鉴权};
    \node[usecase] (project) at (3.75, 1.55) {创建与\\维护项目};
    \node[usecase] (asset) at (3.75, 0.25) {管理图像组\\与图像};
    \node[usecase] (start) at (3.75, -1.05) {启动 Agent\\智能检测};
    \node[usecase] (review) at (3.75, -2.35) {复核\\检测结果};
    \node[usecase] (report) at (3.75, -3.65) {生成与查看\\评估报告};

    \node[usecase] (route) at (9.75, 2.25) {自动选择\\检测任务};
    \node[usecase] (detect) at (9.75, 0.95) {执行原子\\检测能力};
    \node[usecase] (summary) at (9.75, -0.35) {生成单图\\检测总结};
    \node[usecase] (context) at (9.75, -1.65) {整合项目\\报告上下文};
    \node[usecase] (retry) at (9.75, -2.95) {失败重试与\\状态恢复};

    \draw[association] (user.east) -- (login.west);
    \draw[association] (user.east) -- (project.west);
    \draw[association] (user.east) -- (asset.west);
    \draw[association] (user.east) -- (start.west);
    \draw[association] (user.east) -- (review.west);
    \draw[association] (user.east) -- (report.west);
    \draw[association] (user.east) -- (retry.west);

    \draw[include] (start.east) -- node[uml label, above, sloped] {$\langle\langle$include$\rangle\rangle$} (route.west);
    \draw[include] (start.east) -- node[uml label, above, sloped] {$\langle\langle$include$\rangle\rangle$} (detect.west);
    \draw[include] (start.east) -- node[uml label, above, sloped] {$\langle\langle$include$\rangle\rangle$} (summary.west);
    \draw[include] (report.east) -- node[uml label, above, sloped] {$\langle\langle$include$\rangle\rangle$} (context.west);
    \draw[include] (retry.west) -- node[uml label, below, sloped] {$\langle\langle$extend$\rangle\rangle$} (start.east);
    \draw[include] (retry.west) -- node[uml label, below, sloped] {$\langle\langle$extend$\rangle\rangle$} (report.east);
  \end{tikzpicture}
  \caption{系统核心 UML 用例图}\label{fig:requirements-core-use-case}
\end{figure}

这些用例基本对应用户从上传资料到获得评估结论的主要操作，也与\cref{sec:introduction-research-problem}中提出的三个研究问题相呼应。图像资产构建主要处理数据如何进入系统、如何按项目和区域组织；Agent 智能检测关注每张图应该执行哪些检测任务，以及这些任务如何被调度；人工复核和报告生成则把模型结果进一步转化为韧性评估表达。与普通后台管理系统相比，本系统的重点不只是增删改查页面，而是围绕长耗时任务组织业务流程。尤其是 Agent 智能检测环节，背后会经过多个服务、多次回调和多种任务状态，因此需要通过状态机和异步编排来维持流程稳定。

需要说明的是，UML 用例图更适合说明参与者能够使用哪些功能，以及这些功能属于系统边界内还是边界外。它并不直接表示用户必须按照怎样的顺序完成操作。实际使用时，本文系统仍按项目基础信息、图像资产构建、Agent 智能检测、人工复核确认和评估报告生成这一流程推进。每个阶段的进入条件和限制将在后续需求分析中继续说明。

\section{功能需求分析}
\label{sec:requirements-functional-requirement-analysis}

功能需求需要服务于完整业务闭环，而不是只覆盖独立页面。根据前述用例，系统功能可以主要划分为项目与权限、图像资产、智能检测、人工复核与报告生成四组需求。其中，项目与图像资产需求为后续检测提供上下文和数据基础；智能检测需求承载 Agent 自主任务调度方法和原子检测能力；复核与报告需求则承载 LLM 分层信息整合策略的最终输出。四组需求之间存在先后依赖，任何一组缺失都会导致“图像采集完成但评估报告难以形成”的现实问题继续存在。

\subsection{项目与权限需求}
\label{subsec:requirements-project-and-permission-requirement}

项目是所有业务数据的归属边界。系统需要支持用户注册、登录、令牌刷新、头像管理、项目创建、项目删除、项目列表查询和项目详情维护。所有项目相关接口都需要进行用户身份校验，保证用户只能访问自己的项目数据。项目基础信息不仅用于页面展示，也会进入项目级报告上下文，因此建筑名称、位置、建设年份、既有问题和评估目标等字段需要能够被长期保存和读取。

项目阶段推进也需要进行条件校验。系统不能允许用户在没有图像资产时进入智能检测阶段，也不能在检测任务未完成时进入人工复核阶段，更不能在报告未成功生成时将项目视为完成。阶段推进本质上是业务状态机的一部分，需要由后端统一判断，而不是由前端根据页面按钮状态自行决定。

\subsection{图像资产需求}
\label{subsec:requirements-image-asset-requirement}

图像资产构建是后续 Agent 检测的基础。系统需要支持图像组管理、图像上传、图像移动、批量删除、批量移动、原图查看、上传状态展示和上传失败处理。图像组不仅是前端分类展示方式，也承担区域语义表达功能，例如“东立面”“裙楼区域”或“玻璃幕墙区域”等分组名称会影响项目报告中的区域化分析。

由于图像数量可能较多，前端需要支持自适应卡片布局、分组折叠、批量选择和实时状态提示。后端需要通过对象存储保存原图和缩略图，通过数据库保存图像元数据、图像状态和分组关系，通过 WebSocket 推送图像状态变化。这样，图像不再是零散文件，而是具备项目归属、区域归属和处理状态的可追踪资产。

\subsection{智能检测需求}
\label{subsec:requirements-intelligent-detection-requirement}

智能检测阶段需要支持项目级启动检测、失败重试、状态查询和结果查询。启动检测时，系统应为所有已上传图像创建检测主任务。检测流程由 Worker 调度器自主推进，用户不需要逐张图像操作。系统需要先调用分类 Agent 生成受控任务代码，再根据任务代码调度金属锈蚀、石材裂缝、石材污渍、玻璃平整度和玻璃爆裂等原子检测能力。

智能检测需求还包括过程可见性。检测不是瞬时操作，用户需要知道当前图像处于分类、检测、总结、成功或失败中的哪一阶段。状态查询用于页面刷新后的状态恢复，WebSocket 用于页面运行时的实时更新。失败重试需要清理旧任务、旧检测产物和缓存，再重新创建任务，避免旧数据影响新一次检测。

\subsection{复核与报告需求}
\label{subsec:requirements-review-and-report-requirement}

复核阶段需要支持读取和保存图像检测任务的复核信息，包括准确性判断和补充评论。人工复核不是对 Agent 和模型能力的否定，而是面向工程评估责任的必要补充。系统需要保留模型原始结果，同时记录人的判断，使项目报告能够体现“自动检测 + 人工确认”的协同过程。

报告阶段需要支持报告状态查询、失败重试、报告结果展示和报告导出。报告生成本身属于后置任务，在用户推进项目阶段后自动触发。报告输入应包含项目基础信息、图像组信息、图像级总结、原子检测结果和人工复核意见，而不是只拼接原始底层指标。这样设计可以降低用户操作复杂度，使报告生成成为自然流程的一部分，也与\cref{sec:foundation-large-language-model-information-integration-and-report-generation}提出的 LLM 分层信息整合思想保持一致。

\section{阶段化需求分析}
\label{sec:requirements-staged-requirement-analysis}

为了使系统符合土木工程领域用户的真实工作习惯，本文将项目工作流划分为五个阶段：项目基础信息、图像资产构建、Agent 智能检测、人工复核确认和评估报告生成。每个阶段既有独立功能，又为后续阶段提供数据。项目基础信息阶段提供建筑名称、地点、年份和描述等上下文；图像资产阶段提供图像和分组；智能检测阶段产生结构化结果和图像总结；人工复核阶段补充人的判断；报告生成阶段形成最终输出。

\begin{figure}[!htb]
  \centering
  \resizebox{\linewidth}{!}{%
    \begin{tikzpicture}[
      font=\footnotesize,
      >=stealth,
      state/.style={draw, rounded corners=2pt, align=center, minimum width=1.45cm, minimum height=0.62cm, line width=0.45pt, inner sep=2pt},
      stage/.style={state, fill=blue!8, minimum width=1.65cm},
      process/.style={state, fill=orange!10},
      success/.style={state, fill=green!10},
      failed/.style={state, fill=red!8},
      panel/.style={draw, rounded corners=2pt, line width=0.45pt, fill=gray!4},
      panel title/.style={font=\scriptsize, anchor=west},
      arrow/.style={->, line width=0.5pt},
      trigger/.style={->, dashed, gray!70, line width=0.45pt},
      exception/.style={->, dashed, line width=0.45pt},
      label/.style={font=\tiny, align=center, inner sep=1pt}
    ]
      \draw[panel] (0, -0.75) rectangle (12.4, 1.10);
      \node[panel title] at (0.2, 0.80) {项目工作流主状态机};
      \node[stage] (p-profile) at (1.05, 0) {项目基础\\信息};
      \node[stage] (p-assets) at (3.11, 0) {图像资产\\构建};
      \node[stage] (p-detection) at (5.17, 0) {Agent\\智能检测};
      \node[stage] (p-review) at (7.23, 0) {人工复核\\确认};
      \node[stage] (p-report) at (9.29, 0) {评估报告\\生成};
      \node[success, minimum width=1.65cm] (p-done) at (11.35, 0) {项目\\完成};

      \draw[arrow] (p-profile) -- (p-assets);
      \draw[arrow] (p-assets) -- (p-detection);
      \draw[arrow] (p-detection) -- (p-review);
      \draw[arrow] (p-review) -- (p-report);
      \draw[arrow] (p-report) -- (p-done);

      \draw[panel] (0, -5.20) rectangle (6.05, -1.05);
      \node[panel title] at (0.2, -1.35) {图像资产构建状态机};
      \node[process] (u-submit) at (1.15, -2.45) {提交\\上传任务};
      \node[process] (u-uploading) at (3.025, -2.45) {上传中};
      \node[success] (u-success) at (4.90, -2.45) {上传成功};
      \node[failed] (u-failed) at (3.025, -3.75) {上传失败};
      \node[process] (u-retry) at (4.90, -3.75) {重新提交\\上传任务};

      \draw[arrow] (u-submit) -- (u-uploading);
      \draw[arrow] (u-uploading) -- (u-success);
      \draw[exception] (u-uploading) -- (u-failed);
      \draw[exception] (u-failed) -- (u-retry);
      \draw[exception] (u-retry.north west) -- (u-uploading.south east);

      \draw[panel] (6.35, -5.20) rectangle (12.40, -1.05);
      \node[panel title] at (6.55, -1.35) {检测主任务状态机};
      \node[process] (m-created) at (7.55, -2.15) {主任务\\创建};
      \node[process] (m-classifying) at (9.375, -2.15) {Agent\\分类中};
      \node[process] (m-subtask-created) at (11.20, -2.15) {子任务\\创建};
      \node[process] (m-detecting) at (11.20, -3.25) {子任务\\检测中};
      \node[process] (m-summarizing) at (9.375, -3.25) {单图\\总结中};
      \node[success] (m-success) at (7.55, -4.35) {主任务\\成功};
      \node[process] (m-retry) at (9.375, -4.35) {失败\\重试};
      \node[failed] (m-failed) at (11.20, -4.35) {主任务\\失败};

      \draw[arrow] (m-created) -- (m-classifying);
      \draw[arrow] (m-classifying) -- (m-subtask-created);
      \draw[arrow] (m-subtask-created) -- (m-detecting);
      \draw[arrow] (m-detecting) -- (m-summarizing);
      \draw[arrow] (m-summarizing) -- (m-success);
      \draw[exception] (m-classifying.south east) -- (m-failed.north west);
      \draw[exception] (m-detecting) -- (m-failed);
      \draw[exception] (m-summarizing.south east) -- (m-failed.north west);
      \draw[exception] (m-failed) -- (m-retry);
      \draw[exception] (m-retry.north west) -- (m-created.south);

      \draw[panel] (0, -8.90) rectangle (6.05, -5.50);
      \node[panel title] at (0.2, -5.80) {检测子任务状态机};
      \node[process] (s-waiting) at (1.15, -6.75) {子任务\\待执行};
      \node[process] (s-running) at (3.025, -6.75) {模型\\推理中};
      \node[process] (s-uploading) at (4.90, -6.75) {产物\\上传};
      \node[process] (s-checking) at (3.025, -8.05) {推理产物\\完整性校验};
      \node[success] (s-success) at (4.90, -8.05) {子任务\\成功};
      \node[failed] (s-failed) at (1.15, -8.05) {子任务\\失败};

      \draw[arrow] (s-waiting) -- (s-running);
      \draw[arrow] (s-running) -- (s-uploading);
      \draw[arrow] (s-uploading.south west) -- (s-checking.north east);
      \draw[arrow] (s-checking) -- (s-success);
      \draw[exception] (s-running.south west) -- (s-failed.north east);
      \draw[exception] (s-checking) -- (s-failed);
      \draw[exception] (s-failed) -- (s-waiting);

      \draw[panel] (6.35, -8.90) rectangle (12.40, -5.50);
      \node[panel title] at (6.55, -5.80) {评估报告生成状态机};
      \node[process] (r-waiting) at (7.55, -6.75) {报告\\待生成};
      \node[process] (r-generating) at (9.375, -6.75) {报告\\生成中};
      \node[success] (r-success) at (11.20, -6.75) {报告生成\\成功};
      \node[failed] (r-failed) at (9.375, -8.05) {报告生成\\失败};
      \node[process] (r-retry) at (11.20, -8.05) {失败\\重试};

      \draw[arrow] (r-waiting) -- (r-generating);
      \draw[arrow] (r-generating) -- (r-success);
      \draw[exception] (r-generating) -- (r-failed);
      \draw[exception] (r-failed) -- (r-retry);
      \draw[exception] (r-retry.north west) -- (r-generating.south east);
    \end{tikzpicture}%
  }
  \caption{项目工作流与任务执行状态机总览}\label{fig:requirements-stage-state-machine}
\end{figure}

如\cref{fig:requirements-stage-state-machine} 所示，阶段化需求不仅包含页面阶段之间的推进关系，还包含每个阶段内部的细粒度状态机。图像资产阶段需要处理上传成功和上传失败后的重新提交；Agent 智能检测阶段需要处理主任务的分类、子任务创建、子任务检测和图像级总结；检测子任务自身还包含模型执行、产物上传和完整性校验；报告生成阶段则需要区分生成中、生成成功和生成失败重试。只有将这些状态机统一纳入后端业务控制，系统才能在长耗时、跨服务和可失败的任务链路中保持一致。

这种阶段化设计的意义在于把复杂评估任务拆分为用户可理解的步骤。用户不需要一次性理解所有数据对象，而是在流程推进中逐步完成输入、确认和输出。对于系统而言，阶段边界也提供了校验条件。例如，没有上传图像时不能进入检测阶段；检测任务未全部成功时不能进入人工复核阶段；只有报告生成成功后，项目才能被视为最终完成。阶段化需求因此同时承担用户体验组织和业务状态约束两个作用。

阶段化设计还使前端交互更清晰。图像资产阶段允许编辑、移动和删除图像；进入检测阶段后，资产编辑能力应逐步收敛；进入复核阶段后，用户关注点从数据准备转为结果确认；进入报告阶段后，用户关注点转为阅读和导出结论。不同阶段的按钮、权限和状态展示需要与业务含义一致。

\section{非功能需求分析}
\label{sec:requirements-nonfunctional-requirement-analysis}

\begin{figure}[!htb]
  \centering
  \resizebox{\linewidth}{!}{%
    \begin{tikzpicture}[
      font=\scriptsize,
      >=stealth,
      panel/.style={draw, rounded corners=2pt, line width=0.45pt, fill=gray!4},
      target/.style={draw, rounded corners=2pt, align=center, minimum width=2.7cm, minimum height=0.86cm, line width=0.45pt, fill=blue!6, inner sep=3pt},
      requirement/.style={draw, rounded corners=2pt, align=center, minimum width=2.45cm, minimum height=0.64cm, line width=0.45pt, fill=green!7, inner sep=2pt},
      constraint/.style={draw, rounded corners=2pt, align=center, minimum width=4.1cm, minimum height=0.64cm, line width=0.45pt, fill=orange!8, inner sep=2pt},
      arrow/.style={->, line width=0.5pt},
      title/.style={font=\scriptsize\bfseries, align=center}
    ]
      \draw[panel] (0.0, 0.55) rectangle (11.85, -4.65);
      \node[title] at (1.9, 0.15) {运行目标};
      \node[title] at (5.375, 0.15) {非功能需求};
      \node[title] at (9.25, 0.15) {工程约束};

      \node[target] (target) at (1.9, -2.35) {真实场景试用\\稳定支撑};

      \node[requirement] (scalability) at (5.375, -0.75) {可扩展性\\能力接入};
      \node[constraint] (scalability-c) at (9.25, -0.75) {任务代码、注册表\\模型与前端可扩展};

      \node[requirement] (consistency) at (5.375, -1.55) {一致性\\状态可信};
      \node[constraint] (consistency-c) at (9.25, -1.55) {回调落库、事务推进\\统一业务状态};

      \node[requirement] (observability) at (5.375, -2.35) {可观测性\\故障定位};
      \node[constraint] (observability-c) at (9.25, -2.35) {Request ID、任务日志\\WebSocket 与 MQ 追踪};

      \node[requirement] (deployability) at (5.375, -3.15) {可部署性\\环境落地};
      \node[constraint] (deployability-c) at (9.25, -3.15) {容器编排、模型挂载\\服务依赖配置};

      \node[requirement] (experience) at (5.375, -3.95) {用户体验\\过程反馈};
      \node[constraint] (experience-c) at (9.25, -3.95) {进度展示、状态恢复\\结果解释与追溯};

      \draw[arrow] (target.east) -- (scalability.west);
      \draw[arrow] (target.east) -- (consistency.west);
      \draw[arrow] (target.east) -- (observability.west);
      \draw[arrow] (target.east) -- (deployability.west);
      \draw[arrow] (target.east) -- (experience.west);

      \draw[arrow] (scalability) -- (scalability-c);
      \draw[arrow] (consistency) -- (consistency-c);
      \draw[arrow] (observability) -- (observability-c);
      \draw[arrow] (deployability) -- (deployability-c);
      \draw[arrow] (experience) -- (experience-c);
    \end{tikzpicture}%
  }
  \caption{非功能需求与工程约束关系}
  \label{fig:requirements-nonfunctional-requirements-overview}
\end{figure}

除业务流程能够走通之外，系统还需要考虑一些更偏工程运行层面的要求，包括可扩展性、一致性、可观测性、可部署性和用户体验等。这里的非功能需求并不是附加项，而是和\cref{sec:foundation-agent-engineering-controllability-problem}中讨论的 Agent 工程化问题直接相关。原因在于，Agent 输出存在一定不确定性，检测和报告生成往往不是瞬时完成，图像原图与检测产物也会占用较多存储空间，同时多个服务之间还存在回调和状态同步。如果只关注页面功能而忽略非功能需求，系统即使能够完成一次演示，也难以稳定支撑真实场景试用。

为便于理解，\cref{fig:requirements-nonfunctional-requirements-overview} 将非功能需求概括为围绕真实场景试用形成的五类工程约束，分别对应能力扩展、状态可信、故障定位、环境落地和过程反馈等问题。下面对五类非功能需求进行详细说明。

\subsection{可扩展性}
\label{subsec:requirements-scalability}

建筑幕墙检测能力具有持续扩展需求。当前系统包含五类原检测子任务，但未来可能新增连接件松动、密封胶老化、幕墙变形、热工异常等检测能力。因此系统不能将检测流程硬编码在某个单体函数中，而应通过任务代码、统一接口和注册表机制组织能力。Agent 输出也应限制在系统已注册能力范围内，以便新增能力时只需补充任务代码、模型服务、数据库字段和前端展示逻辑。可扩展性是系统持续接入新检测能力并适应真实工程场景演进的重要基础。

\subsection{一致性}
\label{subsec:requirements-consistency}

检测任务涉及多服务协作和异步回调。一旦状态更新顺序混乱，前端进度和数据库结果就可能不一致。本文采用“先建记录、后调服务、回调落库、事务推进”的原则。即在调用外部 Activity 服务前，BIZ 服务先创建任务记录并设置状态；Activity 服务只通过回调上报结果；BIZ 在事务中更新结果并判断是否推进主任务。该设计保证了系统状态由业务中心统一掌握。一致性是多服务异步协作场景下业务状态可信推进的重要前提。

\subsection{可观测性}
\label{subsec:requirements-observability}

Agent 和模型调用存在外部依赖，失败原因可能来自网络、模型、对象存储、数据库或第三方 Agent 平台。因此系统需要完整日志和状态追踪。每个请求需要携带 Request ID 并在 HTTP、RPC 和回调链路中透传；每个任务需要记录 UUID、状态、开始时间、完成时间和运行耗时；每次 WebSocket 广播和 RocketMQ 事件也需要有日志。可观测性是长链路系统稳定运行的重要保障。

\subsection{可部署性}
\label{subsec:requirements-deployability}

毕业设计不是脚本演示，系统需要能够在服务器上通过容器方式启动。由于系统由多个服务组成，部署需要处理镜像构建、服务依赖、环境变量、消息队列 Topic 初始化、模型文件挂载和网络端口映射。尤其是 Reasoning 服务同时依赖 Golang 运行时、Python 环境和模型权重，部署方案需要兼顾镜像体积和模型文件管理。本文采用服务镜像加模型挂载的方式，使模型文件可以独立维护，避免镜像过大。

\subsection{用户体验}
\label{subsec:requirements-user-experience}

用户体验需求主要体现在等待过程和结果解释上。检测任务耗时不可避免，因此前端需要以进度条、流程节点和图像蒙层展示状态。结果解释方面，用户需要同时看到原图、检测产物、结构化指标和 Agent 总结。对于区域类结果，系统需要支持在总体数据和各区域数据之间切换，使用户能够从总体结论追踪到具体区域。

\section{关键挑战}
\label{sec:requirements-key-challenges}

根据上述需求分析中可以看出，本文系统不能简单理解为普通 CRUD 页面组合。它更接近一个围绕图像资产、异步任务和报告生成组织起来的业务流程系统，因此实现难点主要集中在以下几个方面。

第一点是业务链路较长。用户上传图像后，系统还需要完成图像归档、分类路由、并发检测、图像总结、人工复核和项目报告生成等步骤。任一环节失败或状态没有及时更新，都会影响后续阶段的推进。

第二个难点是服务协作较多。系统中前端负责交互，API 负责接入和鉴权，BIZ 负责业务状态，三个 Activity 服务分别处理分类、推理和总结，数据库、对象存储、缓存和消息队列又承担不同的数据支撑职责。各模块边界不同，但需要在一次检测流程中连续配合。

第三，状态同步要求较高。检测状态既要持久化到数据库，又要通过 RocketMQ 触发事件，并通过 WebSocket 推送到前端。用户刷新页面后仍应看到正确进度，因此系统不能只依赖前端内存状态，而要以后端任务状态作为最终依据。

第四，结果结构并不统一。五类检测任务的输出字段、区域数量和产物图像数量都不完全相同，有的输出掩码和面积占比，有的输出分类置信度或校正图。报告生成阶段还需要把这些差异化结果重新整理成统一上下文，否则后续 LLM 很难稳定生成项目级结论。

第五，报告输入组织复杂。项目级报告并不是简单拼接所有检测结果，而是需要把原子指标、图像总结、图像组信息和人工复核意见组织为模型可理解的上下文。上下文过短会导致报告缺少依据，上下文过长又可能造成重点丢失。

第六，跨语言能力集成复杂。视觉检测模型主要运行在 Python 环境中，而后端业务服务采用 Golang 实现。系统需要让 Python 专注模型推理，让 Golang 负责下载、上传、回调和状态管理。跨语言边界既要简洁，又要保证错误能够被 Golang 层捕获并反馈给上游状态机。

因此，本文的工作量主要体现在跨模块的工程链路设计和 Agent 能力嵌入。系统需要同时满足“智能性”和“工程稳定性”：Agent 负责语义判断和信息整合，工程系统负责协议、状态、存储和容错。二者缺一不可。上述挑战也说明，后续技术方案不能只讨论 Agent 提示词，也不能只讨论工程接口，而必须把 Agent 方法与工程状态机共同设计。

\section{数据流分析}
\label{sec:requirements-data-flow-analysis}

数据流分析用于说明系统中的数据如何从用户输入转化为可复核的检测结果和可阅读的评估报告。与普通信息管理系统不同，本文系统同时包含前端交互数据、业务状态数据、图像文件数据、模型推理结果、Agent 语义输出和实时事件数据。为了体现这些数据之间的关系，\cref{fig:requirements-system-data-flow} 从 FE 前端层、API 接入层、BIZ 业务层、Activity 基础能力服务和基础设施五个层次刻画系统主要数据流。

如\cref{fig:requirements-system-data-flow} 所示，系统数据流首先从用户在前端发起项目操作和图像上传开始。前端通过 Core API 获取业务数据和上传授权，并将大体量图像文件直接写入 MinIO；Core API 负责鉴权、HTTP 接入和 WebSocket 连接维护，复杂业务状态则下沉到 Core BIZ。Core BIZ 在 MySQL 中保存项目、图像、检测主任务、检测子任务、复核记录和报告状态，在 Redis 中配合完成令牌状态和临时控制，并根据业务阶段创建分类、检测、总结和报告任务。

\begin{figure}[!htb]
  \centering
  \resizebox{\linewidth}{!}{%
    \begin{tikzpicture}[
      font=\footnotesize,
      >=stealth,
      role/.style={draw, rounded corners=2pt, align=center, minimum width=2.6cm, minimum height=0.95cm, line width=0.5pt, inner sep=2pt},
      service/.style={draw, rounded corners=2pt, align=center, minimum width=2.75cm, minimum height=0.95cm, line width=0.5pt, inner sep=2pt},
      storage/.style={draw, rounded corners=2pt, align=center, minimum width=2.55cm, minimum height=0.88cm, line width=0.5pt, inner sep=2pt},
      arrow/.style={->, line width=0.55pt},
      callback/.style={->, line width=0.58pt, draw=red!70!black},
      object/.style={<->, dashed, line width=0.58pt, draw=blue!70!black},
      object read/.style={->, dashed, line width=0.58pt, draw=blue!70!black},
      mq event/.style={->, dashed, line width=0.58pt, draw=orange!85!black},
      ws event/.style={->, dashed, line width=0.58pt, draw=cyan!70!black},
      state/.style={<->, dashed, line width=0.58pt, draw=green!45!black},
      note/.style={align=center, font=\footnotesize},
      edge label/.style={font=\tiny, inner sep=1pt},
      legend/.style={font=\tiny, align=left}
    ]
      \node[role, fill=blue!8] (user) at (0, 3) {FE 前端\\项目创建、图像上传\\结果查看与人工复核};
      \node[service, fill=cyan!8] (api) at (3.85, 3) {Core API\\HTTP 接入\\鉴权与 WebSocket};
      \node[service, fill=green!8, minimum width=3.1cm, minimum height=1.2cm] (biz) at (7.75, 3) {Core BIZ\\业务状态机\\任务创建与事务落库};

      \node[service, fill=orange!10] (classification) at (12.0, 5.25) {Classification\\分类 Agent\\输出检测任务代码};
      \node[service, fill=red!7] (reasoning) at (12.0, 3) {Reasoning\\检测模型服务\\执行原子检测};
      \node[service, fill=purple!8] (summary) at (12.0, 0.75) {Summary\\总结 Agent\\图像级总结与项目报告};

      \node[storage, fill=pink!14, minimum width=2.35cm] (websocket) at (1.92, 1.40) {WebSocket\\实时状态通道};
      \node[storage, fill=gray!10, anchor=west] (mysql) at (-0.80, -1.25) {MySQL\\项目 / 图像 / 任务\\指标 / 复核 / 报告};
      \path (mysql.east) ++(0.78, 0) node[storage, fill=gray!10, anchor=west] (redis) {Redis\\令牌状态\\缓存与临时控制};
      \path (redis.east) ++(0.78, 0) node[storage, fill=gray!10, anchor=west] (minio) {MinIO\\原图 / 缩略图 / 检测产物\\结构化多源分析数据};
      \path (minio.east) ++(0.78, 0) node[storage, fill=yellow!14, anchor=west] (rocketmq) {RocketMQ\\项目事件 Topic};

      \draw[arrow] ([yshift=0.16cm]user.east) -- node[edge label, above=2pt] {操作请求} ([yshift=0.16cm]api.west);
      \draw[arrow] ([yshift=-0.16cm]api.west) -- node[edge label, below=2pt] {状态展示} ([yshift=-0.16cm]user.east);
      \draw[arrow] ([yshift=0.16cm]api.east) -- node[edge label, above=2pt] {RPC 调用} ([yshift=0.16cm]biz.west);
      \draw[arrow] ([yshift=-0.16cm]biz.west) -- node[edge label, below=2pt] {查询结果} ([yshift=-0.16cm]api.east);

      \draw[arrow] ([xshift=-0.16cm]biz.north east) -- node[edge label, pos=0.55, above left=1pt] {图像路由任务} ([yshift=0.16cm]classification.south west);
      \draw[callback] ([xshift=0.16cm]classification.south west) -- node[edge label, pos=0.55, below right=1pt] {任务代码} ([yshift=-0.16cm]biz.north east);

      \draw[arrow] ([yshift=0.16cm]biz.east) -- node[edge label, above=2pt] {检测子任务} ([yshift=0.16cm]reasoning.west);
      \draw[callback] ([yshift=-0.16cm]reasoning.west) -- node[edge label, below=2pt] {指标与产物} ([yshift=-0.16cm]biz.east);

      \draw[callback] ([xshift=0.16cm]summary.north west) -- node[edge label, pos=0.55, above right=1pt] {总结与报告} ([yshift=0.16cm]biz.south east);
      \draw[arrow] ([xshift=-0.16cm]biz.south east) -- node[edge label, pos=0.55, below left=1pt] {总结/报告任务} ([yshift=-0.16cm]summary.north west);

      \draw[state] ([xshift=-0.36cm]biz.south) -- node[edge label, left=2pt] {状态读写} (mysql.north east);
      \draw[object] (minio.north) -- node[edge label, right=2pt] {文件读写} ([xshift=0.18cm]biz.south);
      \draw[object] (user.south east) -- (minio.north west);
      \draw[object read] ([xshift=-0.22cm]minio.north east) -- (classification.west);
      \draw[object] (minio.north east) -- (reasoning.west);
      \draw[object] ([yshift=-0.22cm]minio.north east) -- (summary.west);
      \draw[state] ([xshift=-0.10cm]biz.south) -- node[edge label, pos=0.55, right=2pt] {缓存读写} (redis.north east);
      \draw[mq event] ([xshift=0.42cm]biz.south) -- (rocketmq.north);
      \draw[mq event] (rocketmq.north west) -- (api.south east);
      \draw[ws event] (api.south) -- node[edge label, right=1pt] {状态广播} (websocket.east);
      \draw[ws event] (websocket.west) -- node[edge label, left=1pt] {实时刷新} (user.south);

      \node[font=\tiny, anchor=center, inner sep=0pt] at (6.00, -2.58) {%
        \begin{tabular}{@{}c@{\hspace{0.10cm}}l@{\hspace{0.38cm}}c@{\hspace{0.10cm}}l@{\hspace{0.38cm}}c@{\hspace{0.10cm}}l@{\hspace{0.38cm}}c@{\hspace{0.10cm}}l@{\hspace{0.38cm}}c@{\hspace{0.10cm}}l@{\hspace{0.38cm}}c@{\hspace{0.10cm}}l@{}}
          \tikz[baseline=-0.55ex]{\draw[->, line width=0.55pt] (0,0) -- (0.34,0);} & HTTP / RPC 控制流 &
          \tikz[baseline=-0.55ex]{\draw[->, line width=0.58pt, draw=red!70!black] (0,0) -- (0.34,0);} & 回调控制流 &
          \tikz[baseline=-0.55ex]{\draw[->, dashed, line width=0.58pt, draw=blue!70!black] (0,0) -- (0.34,0);} & 对象存储文件流 &
          \tikz[baseline=-0.55ex]{\draw[->, dashed, line width=0.58pt, draw=orange!85!black] (0,0) -- (0.34,0);} & MQ 事件流 &
          \tikz[baseline=-0.55ex]{\draw[->, dashed, line width=0.58pt, draw=cyan!70!black] (0,0) -- (0.34,0);} & WebSocket 事件流 &
          \tikz[baseline=-0.55ex]{\draw[->, dashed, line width=0.58pt, draw=green!45!black] (0,0) -- (0.34,0);} & 状态与缓存读写
        \end{tabular}%
      };
    \end{tikzpicture}%
  }
  \caption{系统核心数据流与事件流图示}\label{fig:requirements-system-data-flow}
\end{figure}

Activity 服务与 BIZ 之间形成受控的异步执行关系。Classification 服务只返回图像对应的任务代码，Reasoning 服务只返回原子检测指标和产物清单，Summary 服务只返回结构化 JSON 总结或项目报告；所有结果都通过回调进入 BIZ，由 BIZ 在事务中解析、落库并推进状态。图像原图、缩略图、检测产物和项目报告源文件等文件数据统一存储在 MinIO，模型服务在执行时读取受控授权下的原图，并在产物上传完整后再回调结构化结果。

数据流中还存在一条独立的实时事件链路。BIZ 在状态变化后向 RocketMQ 发布项目事件，Core API 消费事件后通过 WebSocket 广播给前端，使用户能够实时看到图像上传、检测进度和报告状态变化。页面刷新或 WebSocket 断连后，前端仍然可以重新调用列表和详情接口，从 MySQL 中恢复完整状态。因此，RocketMQ 和 WebSocket 主要承担增量通知职责，数据库中的结构化状态仍然是系统最终可信来源。

从该数据流可以看出，系统中并不是只有数据库记录在流动，还同时存在文件数据、结构化数据和语义数据三类内容。原图、缩略图、检测产物和报告源文件属于文件数据，主要存储在 MinIO 中，后续复核和结果展示都需要依赖这些文件。项目、图像、任务状态、指标字段、区域 JSON 和复核意见属于结构化数据，保存在 MySQL 中，用来支撑页面查询、状态追踪和流程推进。图像级总结和项目评估报告则属于语义数据，由 Agent 生成后再经过 BIZ 校验并写入数据库。三类数据的作用并不相同：文件数据保留视觉证据，结构化数据记录过程和指标，语义数据负责把底层结果转化为面向运维决策的表达。

\section{约束条件}
\label{sec:requirements-constraints}

系统设计还需要满足若干约束。第一，模型服务和业务服务应保持边界清晰。Python 模型不应直接操作数据库和对象存储业务 Key，否则后续新增模型会带来大量耦合。第二，Agent 输出必须可校验。分类结果需要解析为任务代码，报告结果需要通过固定接口回调，不能让自由文本直接控制状态。第三，检测任务需要支持失败状态和重试扩展。第四，前端状态需要能从服务端恢复，不能只依赖 WebSocket 增量消息，因为用户可能刷新页面或断线重连。第五，报告生成必须以系统已有数据为依据，不能让模型脱离项目上下文自由编造工程背景和检测结论。

\section{本章小结}
\label{sec:requirements-chapter-summary}

本章在摘要、绪论和相关研究基础之上，将“项目级图像资产管理、Agent 自主任务调度、LLM 分层信息整合”三条主线转化为系统需求。通过业务对象和用户流程建模，本章明确了系统角色责任、典型用例、功能需求、非功能需求、数据流和约束条件。需求分析表明，本文系统的核心难点不只是提供图像上传和结果展示功能，而是需要在多阶段业务流程中实现 Agent 自主任务调度方法、原子检测、人工复核和报告生成的闭环协同。

\begin{table}[!htb]
  \caption{需求与设计目标映射}\label{tab:requirement-goal-map}
  \centering
  \begin{tabular}{@{}p{0.14\linewidth}p{0.3\linewidth}p{0.5\linewidth}@{}} \toprule
    \textbf{需求类型} & \textbf{现实问题} & \textbf{设计目标} \\ \midrule
    图像资产管理 & 图像数量多、组织混乱、难以复核 & 建立项目、图像组和图像三级资产结构 \\
    检测任务调度 & 不同材质需要不同检测能力 & 使用分类 Agent 输出受控任务代码 \\
    Agent 可控性 & 模型输出存在不确定性 & 通过工具集合、结构化输出和 BIZ 状态机限制业务后果 \\
    长任务执行 & 检测和总结耗时较长 & 使用 Worker、异步启动和回调推进状态 \\
    多源结果整合 & 底层指标难以直接形成报告 & 采用图像级总结到项目报告的 LLM 分层信息整合策略 \\
    用户感知 & 长时间执行容易造成等待焦虑 & 使用 WebSocket、进度条和流程节点实时展示 \\
    可追踪性 & 多服务链路失败难以定位 & 使用 UUID、状态机、日志和数据库记录追踪全流程 \\ \bottomrule
  \end{tabular}
\end{table}

如\cref{tab:requirement-goal-map} 所示，本章并不是把系统功能简单列成清单，而是把前文提到的现实问题逐一对应到后续设计中。项目级管理、任务调度、长任务执行、多源结果整合、用户感知和可追踪性，都会在系统中落到具体的工程或 Agent 设计点上。因此，\cref{chap:design}将继续讨论面向建筑幕墙多材质图像的 Agent 自主任务调度方法和 LLM 分层信息整合策略，\cref{chap:implementation}则说明这些需求如何进一步落实到服务架构、协议、数据库状态机、对象存储和前端可视化实现中。
